% Provenance (never printed): see sources/outline_sources.md and sources/book1/BRIEFS.md (g6/02).
\chapter{Variables and Input}\label{ch:g6:variables-and-input}

Noor is playing a game on the tablet. At the top of the screen, a word and a
number: \emph{score 0}. She catches a star: \emph{score 10}. Another one:
\emph{score 20}. The number keeps changing, but its name stays the same, and the
game always knows where to find it. \omterm{def:g4:programs-and-operating-system:program}{Programs} keep their numbers and words in
the same way, under names.

\begin{recall}
\code{print} shows text and the results of calculations; Python knows
\code{+}, \code{-}, \code{*}, \code{/}, \code{//} and \texttt{\%}; the lines of a
\omterm{def:g4:programs-and-operating-system:program}{program} run from top to bottom; an \omterm{def:g6:first-python-program:error}{error message} names the line and the kind
of problem (\cref{ch:g6:first-python-program}).
\end{recall}

\section{A name for a value}

\begin{definition}[Value and variable]\label{def:g6:variables-and-input:value}
A \emph{value}\index{value} is a piece of data a \omterm{def:g4:programs-and-operating-system:program}{program} works with: a number
such as \code{25} or \code{2.5}, or a text such as \code{"Noor"}. A
\emph{variable}\index{variable} is a name that a \omterm{def:g4:programs-and-operating-system:program}{program} gives to a value, to
keep it and use it later.
\end{definition}

\begin{definition}[Assignment]\label{def:g6:variables-and-input:assignment}
An \emph{assignment}\index{assignment} is an \omterm{def:g1:instructions-in-order:instruction}{instruction} that gives a \omterm{def:g6:variables-and-input:value}{value} to
a \omterm{def:g6:variables-and-input:value}{variable}. It is written with the sign \code{=}: the name on the left, the
\omterm{def:g6:variables-and-input:value}{value} on the right, as in \code{score = 0}.
\end{definition}

\Needspace*{8\baselineskip}
\omcode{code/grade-6/02-variables-and-input/listings/score.py}{1}{5}{Two
\omterm{def:g6:variables-and-input:value}{variables}.}\label{lst:g6:variables-and-input:score}

\Needspace*{7\baselineskip}
\omoutput{code/grade-6/02-variables-and-input/outputs/score.txt}{Wherever a
\omterm{def:g6:variables-and-input:value}{variable}'s name is written, Python uses its
\omterm{def:g6:variables-and-input:value}{value}.}\label{out:g6:variables-and-input:score}

A \omterm{def:g6:variables-and-input:value}{variable} is like a box with a label: the label is the name, and the box holds
one \omterm{def:g6:variables-and-input:value}{value}. A name is made of letters, digits and the sign \code{_}, starts with
a letter, and is best chosen so that it says what the \omterm{def:g6:variables-and-input:value}{value} means:
\code{score} is better than \code{s}. Capital and small letters make different
names: \code{Score} is not \code{score}.

\section{Assignment is not equality}

In mathematics, $x = x + 10$ is never true. In Python, \code{score = score +
10} is a perfectly good \omterm{def:g1:instructions-in-order:instruction}{instruction}, because \code{=} does not say that two
things are equal: it says \emph{put into}. Python first computes the right side,
with the \omterm{def:g6:variables-and-input:value}{value} the \omterm{def:g6:variables-and-input:value}{variable} has now, then puts the result in the box. The old
\omterm{def:g6:variables-and-input:value}{value} is gone.

\Needspace*{8\baselineskip}
\omcode{code/grade-6/02-variables-and-input/listings/update.py}{1}{5}{Adding 10
three times.}\label{lst:g6:variables-and-input:update}

\Needspace*{5\baselineskip}
\omoutput{code/grade-6/02-variables-and-input/outputs/update.txt}{The box held 0,
then 10, 20 and 30.}\label{out:g6:variables-and-input:update}

\begin{omfigure}
\begin{tikzpicture}[x=1cm, y=1cm, font=\footnotesize,
    vbox/.style={draw=black!70, fill=ocYellow!25, minimum width=1.5cm, minimum height=0.9cm, font=\ttfamily}]
  \node[vbox] (b1) at (0,0) {20};
  \node[above=1pt of b1, font=\ttfamily\footnotesize] {score};
  \node[font=\small] at (0,-0.9) {before};
  \node[draw=black!60, rounded corners=3pt, fill=ocBlue!10, align=center, minimum height=0.9cm] (c) at (3.6,0) {1.\ read the box: 20\\2.\ add 10: 30};
  \node[vbox] (b2) at (7.2,0) {30};
  \node[above=1pt of b2, font=\ttfamily\footnotesize] {score};
  \node[font=\small] at (7.2,-0.9) {after};
  \draw[-{Stealth[length=2.5mm]}, thick] (b1.east) -- (c.west);
  \draw[-{Stealth[length=2.5mm]}, thick] (c.east) -- node[above] {3.\ write} (b2.west);
  \node[anchor=west, font=\ttfamily\small] at (1.0,1.35) {score = score + 10};
  \node[anchor=west, font=\scriptsize, black!60] at (1.0,0.95) {right side first, then the box on the left};
\end{tikzpicture}

{\small The \omterm{def:g6:variables-and-input:assignment}{assignment} \code{score = score + 10}, step by step: the \omterm{def:g6:variables-and-input:value}{value} in
the box is read, 10 is added, and the result replaces the old \omterm{def:g6:variables-and-input:value}{value}.}
\end{omfigure}

To exchange the \omterm{def:g6:variables-and-input:value}{values} of two \omterm{def:g6:variables-and-input:value}{variables}, a third box is needed to keep a copy,
just as two glasses of juice cannot be swapped without a third glass.

\Needspace*{9\baselineskip}
\omcode{code/grade-6/02-variables-and-input/listings/swap.py}{1}{6}{Exchanging
two \omterm{def:g6:variables-and-input:value}{values} with a spare \omterm{def:g6:variables-and-input:value}{variable}.}\label{lst:g6:variables-and-input:swap}

\Needspace*{5\baselineskip}
\omoutput{code/grade-6/02-variables-and-input/outputs/swap.txt}{The two \omterm{def:g6:variables-and-input:value}{values}
have changed places.}\label{out:g6:variables-and-input:swap}

\section{Types}

Not all \omterm{def:g6:variables-and-input:value}{values} are of the same kind. Python tells them apart.

\begin{definition}[Type]\label{def:g6:variables-and-input:type}
The \emph{type}\index{type (of a value)} of a \omterm{def:g6:variables-and-input:value}{value} is the kind of \omterm{def:g6:variables-and-input:value}{value} it
is. Whole numbers have the type \code{int} (for \emph{integer}), decimal
numbers the type \code{float}, and texts in quotes the type \code{str}. The
\omterm{def:g1:instructions-in-order:instruction}{instruction} \code{type(...)} gives the type of a \omterm{def:g6:variables-and-input:value}{value}.
\end{definition}

\Needspace*{7\baselineskip}
\omcode{code/grade-6/02-variables-and-input/listings/kinds.py}{2}{5}{Asking
Python for \omterm{def:g6:variables-and-input:type}{types}.}\label{lst:g6:variables-and-input:kinds}

\Needspace*{8\baselineskip}
\omoutput{code/grade-6/02-variables-and-input/outputs/kinds.txt}{\code{"12"} in
quotes is a text, not a number.}\label{out:g6:variables-and-input:kinds}

\begin{omfigure}
\begin{tikzpicture}[x=1cm, y=1cm, font=\footnotesize]
  \foreach \x/\t/\v/\c/\d in {0/int/42/ocBlue/whole numbers, 4.2/float/2.5/ocMint/decimal numbers, 8.4/str/\code{"Noor"}/ocRed/texts in quotes} {
    \draw[thick, draw=\c!70!black, fill=\c!12, rounded corners=4pt] (\x,0) rectangle ++(3.2,1.3);
    \node[font=\ttfamily\large] at (\x+1.6,0.65) {\v};
    \node[font=\ttfamily\bfseries, text=\c!70!black] at (\x+1.6,1.6) {\t};
    \node at (\x+1.6,-0.3) {\d};
  }
\end{tikzpicture}

{\small Three \omterm{def:g6:variables-and-input:type}{types}, three kinds of box. The same symbol \code{+} does
different things with them.}
\end{omfigure}

The \omterm{def:g6:variables-and-input:type}{type} decides what the operations do. Between two numbers, \code{+} adds;
between two texts, it glues them end to end.

\Needspace*{6\baselineskip}
\omcode{code/grade-6/02-variables-and-input/listings/glue.py}{1}{3}{The same
sign, two \omterm{def:g6:variables-and-input:type}{types}.}\label{lst:g6:variables-and-input:glue}

\Needspace*{7\baselineskip}
\omoutput{code/grade-6/02-variables-and-input/outputs/glue.txt}{3 and 4 make 7;
the texts \code{"3"} and \code{"4"} make the text
\code{"34"}.}\label{out:g6:variables-and-input:glue}

\section{Input and conversions}

The \omterm{def:g1:instructions-in-order:instruction}{instruction} \code{input(...)} prints a question, waits until the user has
typed an answer and pressed \key{Enter}, and gives the answer as a \omterm{def:g6:variables-and-input:value}{value}. In
the \omterm{def:g1:what-is-a-computer:input}{outputs} of this book, what the user typed appears right after the question.

\Needspace*{6\baselineskip}
\omcode{code/grade-6/02-variables-and-input/listings/greet.py}{1}{3}{A \omterm{def:g4:programs-and-operating-system:program}{program}
that asks a question.}\label{lst:g6:variables-and-input:greet}

\Needspace*{7\baselineskip}
\omoutput{code/grade-6/02-variables-and-input/outputs/greet.txt}{Noor typed her
name.}\label{out:g6:variables-and-input:greet}

Line 3 shows a handy way to build a sentence: a text with the letter \code{f}
just before its quotes is an \emph{f-string}, and each name between curly
brackets \texttt{\{\}} is replaced by its \omterm{def:g6:variables-and-input:value}{value}.

\begin{omfigure}
\begin{tikzpicture}[x=1cm, y=1cm, font=\footnotesize,
    flow/.style={draw=black!65, rounded corners=3pt, fill=#1, minimum height=0.9cm, align=center, inner sep=3pt}]
  \node[flow=black!8] (k) at (0,0) {keyboard\\\ttfamily 12};
  \node[flow=ocYellow!30] (i) at (2.3,0) {\ttfamily input};
  \node[flow=ocRed!12] (t) at (4.4,0) {text\\\code{"12"}};
  \node[flow=ocYellow!30] (n) at (6.4,0) {\ttfamily int};
  \node[flow=ocBlue!12] (v) at (8.4,0) {number\\\ttfamily 12};
  \node[flow=ocMint!15] (c) at (10.6,0) {calculation\\\ttfamily 12 + 30};
  \node[flow=black!8] (p) at (12.9,0) {screen\\\ttfamily 42};
  \foreach \a/\b in {k/i, i/t, t/n, n/v, v/c, c/p} \draw[-{Stealth[length=2mm]}, thick] (\a) -- (\b);
\end{tikzpicture}

{\small What the user \omterm{def:g6:variables-and-input:type}{types} reaches the \omterm{def:g4:programs-and-operating-system:program}{program} as text. To compute with it,
the text must first be converted into a number.}
\end{omfigure}

\code{input} always gives a text, even when the user \omterm{def:g6:variables-and-input:type}{types} digits. Adding two
typed numbers then glues them:

\Needspace*{6\baselineskip}
\omcode{code/grade-6/02-variables-and-input/listings/add_wrong.py}{1}{3}{Adding
two typed numbers?}\label{lst:g6:variables-and-input:add-wrong}

\Needspace*{7\baselineskip}
\omoutput{code/grade-6/02-variables-and-input/outputs/add_wrong.txt}{12 and 30
make 1230: two texts were glued.}\label{out:g6:variables-and-input:add-wrong}

\begin{definition}[Type conversion]\label{def:g6:variables-and-input:conversion}
A \emph{type conversion}\index{type conversion} makes, from a \omterm{def:g6:variables-and-input:value}{value}, the \omterm{def:g6:variables-and-input:value}{value}
of another \omterm{def:g6:variables-and-input:type}{type} that stands for the same thing: \code{int("12")} gives the
whole number 12, \code{float("2.5")} the decimal number 2.5, and
\code{str(12)} the text \code{"12"}.
\end{definition}

\Needspace*{6\baselineskip}
\omcode{code/grade-6/02-variables-and-input/listings/add_right.py}{1}{3}{The
typed texts are converted into whole
numbers.}\label{lst:g6:variables-and-input:add-right}

\Needspace*{7\baselineskip}
\omoutput{code/grade-6/02-variables-and-input/outputs/add_right.txt}{Now 12 and
30 make 42.}\label{out:g6:variables-and-input:add-right}

A \omterm{def:g4:programs-and-operating-system:program}{program} that asks, computes and answers is now within reach:

\Needspace*{10\baselineskip}
\omcode{code/grade-6/02-variables-and-input/listings/rectangle.py}{1}{7}{The area
and the perimeter of a typed
rectangle.}\label{lst:g6:variables-and-input:rectangle}

\Needspace*{8\baselineskip}
\omoutput{code/grade-6/02-variables-and-input/outputs/rectangle.txt}{A rectangle
of 7.5 metres by 4 metres.}\label{out:g6:variables-and-input:rectangle}

\section{Following a program by hand}

To understand a \omterm{def:g4:programs-and-operating-system:program}{program}, or to find a mistake in it, we can run it by hand,
like the \omterm{def:g6:first-python-program:interpreter}{interpreter}, writing down the \omterm{def:g6:variables-and-input:value}{value} of each \omterm{def:g6:variables-and-input:value}{variable} after each line.

\begin{definition}[Trace table]\label{def:g6:variables-and-input:trace-table}
A \emph{trace table}\index{trace table} is a table that follows a \omterm{def:g4:programs-and-operating-system:program}{program} by
hand: one column for each \omterm{def:g6:variables-and-input:value}{variable}, one row for each line run, holding the
\omterm{def:g6:variables-and-input:value}{values} of the \omterm{def:g6:variables-and-input:value}{variables} after that line.
\end{definition}

\Needspace*{8\baselineskip}
\omcode{code/grade-6/02-variables-and-input/listings/trace_me.py}{1}{5}{A
\omterm{def:g4:programs-and-operating-system:program}{program} to follow by hand.}\label{lst:g6:variables-and-input:trace-me}

\begin{example}[Tracing five lines]\label{ex:g6:variables-and-input:trace}
A dash means that the \omterm{def:g6:variables-and-input:value}{variable} has no \omterm{def:g6:variables-and-input:value}{value} yet.

\begin{center}
\begin{tabular}{|c|l|c|c|}
\hline
line & \omterm{def:g1:instructions-in-order:instruction}{instruction} & \code{a} & \code{b} \\ \hline
1 & \code{a = 5} & 5 & -- \\
2 & \code{b = a + 3} & 5 & 8 \\
3 & \code{a = b * 2} & 16 & 8 \\
4 & \code{b = b - a} & 16 & $-8$ \\
5 & \code{print(a, b)} & 16 & $-8$ \\ \hline
\end{tabular}
\end{center}

The \omterm{def:g4:programs-and-operating-system:program}{program} should print \code{16 -8}. Running it confirms the table:
\end{example}

\Needspace*{5\baselineskip}
\omoutput{code/grade-6/02-variables-and-input/outputs/trace_me.txt}{The same
\omterm{def:g6:variables-and-input:value}{values} as in the \omterm{def:g6:variables-and-input:trace-table}{trace table}.}\label{out:g6:variables-and-input:trace-me}

\begin{method}[Tracing a program]\label{met:g6:variables-and-input:trace}
Draw one column per \omterm{def:g6:variables-and-input:value}{variable}. Go through the \omterm{def:g4:programs-and-operating-system:program}{program} line by line, in the
order Python runs it; for each line, write a new row: copy the \omterm{def:g6:variables-and-input:value}{values} that do
not change, and write the new \omterm{def:g6:variables-and-input:value}{value} of the \omterm{def:g6:variables-and-input:value}{variable} that is assigned. Compute
the right side of an \omterm{def:g6:variables-and-input:assignment}{assignment} with the \omterm{def:g6:variables-and-input:value}{values} of the row above.
\end{method}

\section{Exercises}

\begin{exercise}[$\star$]\label{exo:g6:variables-and-input:1}
What does this \omterm{def:g4:programs-and-operating-system:program}{program} print?

\omcode{code/grade-6/02-variables-and-input/exercises/three_a.py}{1}{4}{Three
\omterm{def:g6:variables-and-input:assignment}{assignments}.}\label{lst:g6:variables-and-input:three-a}
\end{exercise}

\begin{exercise}[$\star$]\label{exo:g6:variables-and-input:2}
What does \code{print(type(4.0))} print? And \code{print(type(4))}?
\end{exercise}

\begin{exercise}[$\star$]\label{exo:g6:variables-and-input:3}
Try \code{print("5" * 3)} and \code{print(5 * 3)}. What do you see?
\end{exercise}

\begin{exercise}[$\star$]\label{exo:g6:variables-and-input:4}
Write a \omterm{def:g4:programs-and-operating-system:program}{program} that asks for a name and an age, then prints a sentence such
as \code{Sami will be 12 next year.}
\end{exercise}

\begin{exercise}[$\star\star$]\label{exo:g6:variables-and-input:5}
Tom's \omterm{def:g4:programs-and-operating-system:program}{program} asks for two numbers and prints their sum, but for 5 and 7 it
prints 57. What did he forget?
\end{exercise}

\begin{exercise}[$\star\star$]\label{exo:g6:variables-and-input:6}
Write a \omterm{def:g4:programs-and-operating-system:program}{program} that asks for a number of minutes and prints it in hours and
minutes: 135 gives \code{2 hours and 15 minutes}.
\end{exercise}

\begin{exercise}[$\star\star$]\label{exo:g6:variables-and-input:7}
Make the \omterm{def:g6:variables-and-input:trace-table}{trace table} of \cref{lst:g6:variables-and-input:swap}, with the
three \omterm{def:g6:variables-and-input:value}{variables} \code{left}, \code{right} and \code{spare}.
\end{exercise}

\begin{exercise}[$\star\star$]\label{exo:g6:variables-and-input:8}
A notebook costs 2.50 euros. Write a \omterm{def:g4:programs-and-operating-system:program}{program} that asks how many notebooks are
bought and prints the price.
\end{exercise}

\begin{exercise}[$\star\star$]\label{exo:g6:variables-and-input:9}
Lina's \omterm{def:g4:programs-and-operating-system:program}{program} starts with \code{size = int(input("Size: "))}. When she \omterm{def:g6:variables-and-input:type}{types}
\code{2.5}, it stops with a \code{ValueError}. Why? How can she fix it?
\end{exercise}

\begin{exercise}[$\star\star\star$]\label{exo:g6:variables-and-input:10}
How many \omterm{def:g6:variables-and-input:assignment}{assignments} does this \omterm{def:g4:programs-and-operating-system:program}{program} carry out? What does it print?

\omcode{code/grade-6/02-variables-and-input/exercises/count_me.py}{1}{6}{Count
the \omterm{def:g6:variables-and-input:assignment}{assignments}.}\label{lst:g6:variables-and-input:count-me}
\end{exercise}

\begin{exercise}[$\star\star\star$]\label{exo:g6:variables-and-input:11}
Write a \omterm{def:g4:programs-and-operating-system:program}{program} that asks for a number of seconds and prints it in days, hours,
minutes and seconds. What does it print for one million seconds?
\end{exercise}

\begin{exercise}[$\star\star\star$]\label{exo:g6:variables-and-input:12}
Tom wants to exchange the \omterm{def:g6:variables-and-input:value}{values} of \code{a} and \code{b} without a spare
\omterm{def:g6:variables-and-input:value}{variable}, so he writes \code{a = b}, then \code{b = a}. Make the \omterm{def:g6:variables-and-input:trace-table}{trace table}
for \code{a = 3} and \code{b = 8}, and explain what went wrong.
\end{exercise}

\section{Problem: The Clock of the Day}

\begin{problem}[{Weekend problem --- seconds, minutes and hours}]\label{pb:g6:variables-and-input:1}
A stopwatch counts only seconds. We want to read its number as a time of day,
written \code{h:m:s}: hours, minutes and seconds.

\medskip\noindent\textbf{Part I --- By hand.}
\begin{enumerate}
  \item How many seconds are there in one minute? In one hour? In one day?
  \item Write 5,000 seconds in hours, minutes and seconds. Show each division
    with its remainder.
  \item Do the same for 50,000 seconds.
\end{enumerate}

\medskip\noindent\textbf{Part II --- Two programs.} The first \omterm{def:g4:programs-and-operating-system:program}{program} turns
seconds into \code{h:m:s}, the second does the opposite.

\omcode{code/grade-6/02-variables-and-input/listings/clock.py}{1}{7}{From seconds
to \code{h:m:s}.}\label{lst:g6:variables-and-input:clock}

\omcode{code/grade-6/02-variables-and-input/listings/back.py}{1}{5}{From
\code{h:m:s} back to seconds.}\label{lst:g6:variables-and-input:back}

\begin{enumerate}[resume]
  \item Make the \omterm{def:g6:variables-and-input:trace-table}{trace table} of \cref{lst:g6:variables-and-input:clock} when
    the user \omterm{def:g6:variables-and-input:type}{types} 50000. Does it agree with your answer to question 3?
  \item What does \cref{lst:g6:variables-and-input:back} print when the user
    \omterm{def:g6:variables-and-input:type}{types} 13, 53 and 20? Why?
  \item What does the first \omterm{def:g4:programs-and-operating-system:program}{program} print for 3,599 seconds? For 3,600?
  \item The first \omterm{def:g4:programs-and-operating-system:program}{program} prints 13:53:20 in the same way it would print a
    time with a minute such as 7, without a 0 in front: \code{13:7:5}. Which
    real time is \code{13:7:5}?
\end{enumerate}

\medskip\noindent\textbf{Part III --- Why it always works.}
\begin{enumerate}[resume]
  \item Explain why the minutes and the seconds printed by the first \omterm{def:g4:programs-and-operating-system:program}{program}
    are always less than 60.
  \item Explain why the second \omterm{def:g4:programs-and-operating-system:program}{program} always gives back the number typed into
    the first one.
  \item Using the times printed for 0 and for the last second of a day, say
    which \omterm{def:g6:variables-and-input:value}{values} \code{h}, \code{m} and \code{s} can take during one day.
  \item How many different \code{h:m:s} times does one day hold? Check your
    answer with the \omterm{def:g4:programs-and-operating-system:program}{program} below, and state the result.
\end{enumerate}

\omcode{code/grade-6/02-variables-and-input/listings/day.py}{1}{3}{Counting the
times of a day.}\label{lst:g6:variables-and-input:day}
\end{problem}
